\begin{annexitem}{annex:p14.evidence.finite-measures}{thm:p14-evidence-finite-measures}{Finite duality, union laws and allocated bounds}
\begin{proof}
Use the finite, nonempty \(W\) of
\cref{def:p14-evidence-uncertainty-data}. Every mass \(p\) gives
\(P_p(\varnothing)=0\), \(P_p(W)=1\), and
\(0\leq P_p(A)\leq P_p(B)\leq1\) when \(A\subseteq B\).
Taking infima or suprema over the nonempty family preserves these bounds,
normalization and monotonicity. No compactness assumption is needed for
their existence as real numbers, because the values lie in \([0,1]\).
The identity \(P_p(A)=1-P_p(A^c)\) for every mass gives
\[
 \sup_{p\in\mathcal C}P_p(A)
 =\sup_{p\in\mathcal C}(1-P_p(A^c))
 =1-\inf_{p\in\mathcal C}P_p(A^c).
\]
If \(A\cap B=\varnothing\), then for every \(p\in\mathcal C\),
\[
 P_p(A\cup B)=P_p(A)+P_p(B)
 \geq\underline P_{\mathcal C}(A)+\underline P_{\mathcal C}(B).
\]
Taking the infimum proves lower superadditivity. Using the two upper
bounds and taking the supremum proves upper subadditivity. This proves
item 1, including empty events.

For possibility, taking the maximum over a union is taking the larger of
the two separate maxima. The declared empty maximum and nonnegativity make
this identity valid when either event is empty. The normalization of
\(\pi\) gives the value one on \(W\), and the value zero on the empty
event is stipulated. Enlargement cannot decrease a maximum. Taking
complements then gives
\[
 \operatorname{Nec}_{\pi}(A\cap B)
 =1-\operatorname{Poss}_{\pi}(A^c\cup B^c)
 =\min\{\operatorname{Nec}_{\pi}(A),\operatorname{Nec}_{\pi}(B)\}.
\]
Necessity is normalized and monotone by the same complement identity.
All values lie in \([0,1]\). This proves item 2.

For the mass-based constructions, \(m(\varnothing)=0\) gives belief zero
at the empty event. The absence of a subset intersecting the empty event
gives plausibility zero there. Every positive-mass subset is nonempty, so
both values on \(W\) sum all positive masses and equal one. Increasing an
event can only enlarge its containment and intersection index sets, proving
monotonicity. Each nonempty \(B\) either lies in \(A^c\) or intersects
\(A\), exclusively. Consequently
\[
 \operatorname{Pl}_m(A)=\sum_{B\subseteq W}m(B)
          -\sum_{B\subseteq A^c}m(B)=1-\operatorname{Bel}_m(A^c).
\]
Every nonempty \(B\subseteq A\) intersects \(A\); thus the belief sum
is a sub-sum of the plausibility sum. The terms present only in the latter
are exactly the subsets meeting both \(A\) and \(A^c\). This proves
item 3 and its precise gap formula.

To prove item 4, fix a listing of the finite set \(W\). For every positive
mass \(m(B)\), a probability concentrated at the first listed member of
the nonempty \(B\) supplies one allocation. This constructs an element
of \(\mathcal C_m\); it invokes no choice of infinitely many witnesses.
For any allocation, its resulting \(p\) is nonnegative, and finite
interchange of the sums gives
\[
 \sum_{w\in W}p(w)
 =\sum_{B:m(B)>0}m(B)\sum_{w\in B}q_B(w)=1.
\]
For a fixed \(A\), the same interchange gives
\[
 P_p(A)=\sum_{B:m(B)>0}m(B)\sum_{w\in A\cap B}q_B(w).
\]
The inner sum is one if \(B\subseteq A\), zero if
\(B\cap A=\varnothing\), and between zero and one otherwise. Hence
\(\operatorname{Bel}_m(A)\leq P_p(A)\leq\operatorname{Pl}_m(A)\).
For the lower bound, allocate each \(B\not\subseteq A\) to its first
listed member outside \(A\); allocate each \(B\subseteq A\) to any
listed member of \(B\). All inner sums are their indicated lower bounds,
so the resulting mass attains \(\operatorname{Bel}_m(A)\).
For the upper bound, allocate every \(B\) intersecting \(A\) to its
first listed member of \(A\cap B\). The other subsets are disjoint from
\(A\) and contribute zero under any allocation. This attains
\(\operatorname{Pl}_m(A)\). Taking infima and suprema proves the two
identities. The allocations may depend on \(A\) and on the bound being
attained; no claim is made that one mass attains every event bound at once.

For item 5, consider any nonempty positive-mass \(C\). Its indicator
contribution to the claimed inequality satisfies
\[
 \mathbf1_{\{C\subseteq A\cup B\}}
 \geq \mathbf1_{\{C\subseteq A\}}+
        \mathbf1_{\{C\subseteq B\}}-
        \mathbf1_{\{C\subseteq A\cap B\}}.
\]
If \(C\) lies in both events, both sides equal one. If it lies in exactly
one, both sides again equal one. If it lies in neither, the right side is
zero and the left side is zero or one. Multiplying by \(m(C)\) and
summing proves the inequality. Finally, when only singleton subsets have
positive mass, write \(p(w)=m(\{w\})\). Both containment and nonempty
intersection select exactly the singletons in \(A\), so belief and
plausibility both equal \(\sum_{w\in A}p(w)\). Its masses normalize
because all nonsingleton and empty-subset masses are zero.
\end{proof}
\end{annexitem}

\begin{annexitem}{annex:p14.evidence.measure-separation}{thm:p14-evidence-measure-separation}{The complete finite comparison and a non-belief lower envelope}
\begin{proof}
The three components of \(p\) sum to one. The other retained vectors
are nonnegative precisely on the displayed parameter interval and sum
to one there. Allocating \(7/10\) within \(\{a,b\}\) can give any
first mass \(t\in[0,7/10]\), the second mass is then \(7/10-t\),
and the singleton \(\{c\}\) receives its fixed \(3/10\).
These are all and only the members of \(\mathcal C_m\), including both
endpoints. In particular \(p\) is its member with \(t=1/2\).

The ordinary probability column is obtained by adding the indicated
components. For the mass-based columns, a proper event containing both
\(a,b\) includes the \(7/10\) subset mass, while an event containing
\(c\) includes the \(3/10\) singleton mass. These two conditions give
the belief column, including the empty event. An event intersecting
\(\{a,b\}\) includes \(7/10\) in the plausibility column, and one
containing \(c\) includes \(3/10\); these give all eight entries.
The envelope identities of \cref{thm:p14-evidence-finite-measures} also
identify them with the infima and suprema of the vectors just parameterized.

Taking maxima of \(1,3/5,1/5\) over the listed events gives the
possibility column. Taking one minus the maximum over each complement
gives the necessity column. For example, the complement of \(\{a,b\}\)
is \(\{c\}\), so necessity is \(1-1/5=4/5\), whereas the
probability of \(\{a,b\}\) is \(1/2+1/5=7/10\). This verifies
each entry and the stated separation.

For the final counterexample, both four-component vectors are probability
masses. Their values on \(A\) are \(2/5,1/2\), on \(B\) they are
again \(2/5,1/2\), on \(A\cup B\) they are \(2/5,1\), and on
\(A\cap B=\{1\}\) they are \(2/5,0\). Taking minima yields the
four displayed lower-envelope values. A mass-based belief function would
have to satisfy
\[
 2/5=\operatorname{Bel}_m(A\cup B)
 \geq\operatorname{Bel}_m(A)+\operatorname{Bel}_m(B)
       -\operatorname{Bel}_m(A\cap B)=4/5,
\]
which is false. Thus this valid lower probability is not a belief function
of the specified mass form. No assertion about an update or evidence
combination is used.
\end{proof}
\end{annexitem}
